%
% polylog.tex -- template for standalon tikz images
%
% (c) 2021 Prof Dr Andreas Müller, OST Ostschweizer Fachhochschule
%
\documentclass[tikz]{standalone}
\usepackage{amsmath}
\usepackage{times}
\usepackage{txfonts}
\usepackage{pgfplots}
\usepackage{csvsimple}
\definecolor{darkgreen}{rgb}{0,0.6,0}
\definecolor{darkred}{rgb}{0.8,0,0}
\usetikzlibrary{arrows,intersections,math}
\begin{document}
\def\skala{1}
\def\h{3}
\def\v{-0.8}
\def\punkt#1#2{ ({(#1)*\h},{-(#2)*\v}) }
\def\q{\qquad}
\def\d{\displaystyle}
\def\g#1{\bgroup\color{darkgreen}#1\egroup}
\def\r#1{\bgroup\color{darkred}#1\egroup}
\begin{tikzpicture}[>=latex,thick,scale=\skala]

\clip (-6.0,-5.25) rectangle (6.0,4.8);

\newboolean{showgrid}
\setboolean{showgrid}{false}
\def\breite{7}
\def\hoehe{6}

% Gitter
\ifthenelse{\boolean{showgrid}}{
\draw[step=0.1,line width=0.1pt,color=gray!20] (-\breite,-\hoehe) grid (\breite, \hoehe);
\draw[step=0.5,line width=0.4pt,color=gray!20] (-\breite,-\hoehe) grid (\breite, \hoehe);
\draw[color=gray!20]                           (-\breite,-\hoehe) grid (\breite, \hoehe);
\fill[color=gray!20] (0,0) circle[radius=0.05,color=gray!20];
}{}

\node at (0,0) {$\renewcommand{\arraycolsep}{2pt}
\renewcommand{\arraystretch}{3.3}
\begin{array}{lclclclclcl}
                                                            & &                            &\q& b_{l+1} &=& \d-\frac{c_l}{l+1}     & & \r{q_{l+1}}   &=& b_{l+1}           \\
\d\int \g{p_l}                                              &=& d_l + c_l \vartheta        &\q& b_l     &=& \d-\frac{c_{l-1}}{l}   & & \r{q_{l}}     &=& d_l + b_l         \\
\d\int \biggl(\g{p_{l-1}} - ld_l\frac{Du}{u}\biggr)         &=& d_{l-1} + c_{l-1}\vartheta &\q& b_{l-1} &=& \d-\frac{c_{l-2}}{l-1} & & \r{q_{l-1}}   &=& d_{l-1} + b_{l-1} \\
\d\int \biggl(\g{p_{l-2}} - (l-1)d_{l-1}\frac{Du}{u}\biggr) &=& d_{l-2} + c_{l-2}\vartheta &\q& b_{l-2} &=& \d-\frac{c_{l-3}}{l-2} & & \r{q_{l-2}}   &=& d_{l-2} + b_{l-2} \\[-12pt]
\hspace*{40pt}\vdots                                        & &                            &\q&         & &                        &\phantom{\longrightarrow}&&&           \\[-15pt]
\d\int \biggl(\g{p_k} - (k+1)d_{k+1}\frac{Du}{u}\biggr)     &=& d_k + c_k\vartheta         &\q& b_k     &=& \d-\frac{c_{k-1}}{k}   & & \r{q_{k}}     &=& d_k + b_k         \\[-12pt]
\hspace*{40pt}\vdots                                        & &                            &\q&         & &                        & &               & &                   \\[-15pt]
\d\int \biggl(\g{p_1} - 2d_2\frac{Du}{u}\biggr)             &=& d_1 + c_1\vartheta         &\q& b_1     &=& \d-\frac{c_0}{2}       & & \r{q_1}       &=& d_1 + b_1         \\
\d\int \biggl(\g{p_0} - d_1\frac{Du}{u}\biggr)              &=& d_0 + c_0\vartheta         &\q&         & &                        & & \r{\bar{q}_0} &=& d_0 + b_0         
\end{array}$};

% q_l+1
\begin{scope}
\draw[->,line width=0.5pt] (1.0,4.25) to[out=-45,in=180] ++(0.5,-0.2) -- ++(2.6,0) to[out=0,in=-135] ++(0.5,0.2);
\end{scope}

% q_l
\begin{scope}
\draw[->,line width=0.5pt] (-1.7,3.25) to[out=45,in=180] ++(0.6,0.3) -- ++(5.1,0) to[out=0,in=135] ++(0.6,-0.3);
\draw[->,line width=0.5pt] (-1.2,3.2) to[out=55,in=180] ++(1.5,1.45) -- ++(1.8,0);
\draw[->,line width=0.5pt] (-2.0,2.85) to[out=-135,in=90] (-4.3,1.9);
\draw[->,line width=0.5pt] (1.0,2.85) to[out=-45,in=180] ++(0.5,-0.2) -- ++(3.2,0) to[out=0,in=-135] ++(0.5,0.2);
\end{scope}

% q_l-1
\begin{scope}
\draw[->,line width=0.5pt] (-1.7,1.85) to[out=45,in=180] ++(0.6,0.3) -- ++(5.1,0) to[out=0,in=135] ++(0.6,-0.3);
\draw[->,line width=0.5pt] (-0.9,1.8) to[out=55,in=180] ++(1.5,1.45) -- ++(1.3,0);
\draw[->,line width=0.5pt] (1.2,1.45) to[out=-45,in=180] ++(0.5,-0.2) -- ++(3.2,0) to[out=0,in=-135] ++(0.5,0.2);
\draw[->,line width=0.5pt] (-2.0,1.40) to[out=-135,in=90] (-3.4,0.5);
\end{scope}

% q_l-2
\begin{scope}
\draw[->,line width=0.5pt] (-1.7,0.45) to[out=45,in=180] ++(0.6,0.3) -- ++(5.1,0) to[out=0,in=135] ++(0.6,-0.3);
\draw[->,line width=0.5pt] (-0.9,0.4) to[out=55,in=180] ++(1.5,1.45) -- ++(1.3,0);
\draw[->,line width=0.5pt] (1.2,0.05) to[out=-45,in=180] ++(0.5,-0.2) -- ++(3.2,0) to[out=0,in=-135] ++(0.5,0.2);
\draw[->,line width=0.5pt,dotted] (-2.0,0.05) to[out=-135,in=90] (-3.6,-1.35);
\end{scope}

% q_k
\begin{scope}
\draw[->,line width=0.5pt] (-1.7,-1.4) to[out=45,in=180] ++(0.6,0.3) -- ++(5.1,0) to[out=0,in=135] ++(0.6,-0.3);
\draw[->,line width=0.5pt,dotted] (-1.1,-1.40) to[out=55,in=180] ++(2.0,1.88) -- ++(1.1,0);
\draw[->,line width=0.5pt] (1.0,-1.80) to[out=-45,in=180] ++(0.5,-0.2) -- ++(3.2,0) to[out=0,in=-135] ++(0.5,0.2);
\draw[->,line width=0.5pt,dotted] (-2.0,-1.80) to[out=-135,in=90] (-4.5,-3.15);
\end{scope}

% q_1
\begin{scope}
\draw[->,line width=0.5pt] (1.0,-3.65) to[out=-45,in=180] ++(0.5,-0.2) -- ++(3.2,0) to[out=0,in=-135] ++(0.5,0.2);
\draw[->,line width=0.5pt,dotted] (-1.1,-3.25) to[out=55,in=180] ++(1.9,1.88) -- ++(1.1,0);
\draw[->,line width=0.5pt] (-1.7,-3.20) to[out=45,in=180] ++(0.6,0.3) -- ++(5.1,0) to[out=0,in=135] ++(0.6,-0.3);
\draw[->,line width=0.5pt] (-2.0,-3.65) to[out=-135,in=90] (-4.6,-4.6);
\end{scope}

% q_0
\begin{scope}
\draw[->,line width=0.5pt] (-1.7,-4.60) to[out=45,in=180] ++(0.6,0.3) -- ++(5.1,0) to[out=0,in=135] ++(0.6,-0.3);
\draw[->,line width=0.5pt] (-1.15,-4.65) to[out=55,in=180] ++(1.5,1.45) -- ++(1.5,0);
\end{scope}

\end{tikzpicture}
\end{document}

